% SOURCE_STATEMENT: "For an irrational flat torus with a point scatterer, does there exist a subsequence of eigenfunctions that scars in momentum space?"
% Preparation and final self-audit: CONTRIBUTING.md prompts C and D; source/audit artifacts retained in the local physics_statements_2026_10_11 batch.
% !TeX program = lualatex
\documentclass[11pt]{article}
\usepackage[margin=25mm]{geometry}
\usepackage{fontspec}
\setmainfont{Latin Modern Roman}
\usepackage{amsmath,amssymb,mathtools}
\usepackage{unicode-math}
\setmathfont{Latin Modern Math}
\usepackage{xurl,hyperref}
\hypersetup{colorlinks=true,urlcolor=blue}
\setcounter{secnumdepth}{0}
\setlength{\parindent}{0pt}
\setlength{\parskip}{5pt}
\setlength{\emergencystretch}{3em}
\newcommand{\sref}[2]{\hyperlink{source:#1:#2}{[S#2]}}
\newcommand{\cataloguescope}[1]{\paragraph{Recorded target.}#1}
\newcommand{\R}{\mathbb R}
\newcommand{\C}{\mathbb C}
\newcommand{\N}{\mathbb N}
\newcommand{\Z}{\mathbb Z}
\newcommand{\Q}{\mathbb Q}
\newcommand{\F}{\mathbb F}
\newcommand{\eps}{\varepsilon}
\begin{document}
\section{Momentum-Space Scarring on Irrational Point-Scatterer Tori}\label{statement}

\subsection{Definitions and mathematical statement}

\paragraph{Shared notation.}$\mathbb N=\{1,2,\ldots\}$; the standard number fields have their usual meanings. All additional conventions are stated locally.


\paragraph{Definitions and conventions.} An irrational rectangular flat torus has a nonrational squared aspect ratio. A point scatterer is a selfadjoint point-interaction extension of the Laplacian at a fixed torus point, with a coupling convention. Eigenfunctions have Fourier coefficients indexed by the dual lattice. Their normalized squared coefficients determine momentum probability measures on directions; a momentum scar is a subsequential limit with concentration beyond the uniform directional distribution, in the sense examined by the source.

\paragraph{Source scope and qualifications.} The primary report states a conditional positive result assuming Poisson spacing of the unperturbed eigenvalues. The unconstrained catalogue question omits the irrationality and coupling details needed for a full theorem; this preparation does not certify that no later unconditional result exists.

\paragraph{Catalogue mathematical target.}
\paragraph{Catalogue wording retained for traceability.}
For an irrational flat torus with a point scatterer, does there exist a subsequence of eigenfunctions that scars in momentum space?

\paragraph{Saved literature context.} The supplied catalogue's assessment dated 2026-08-21 records: The record is incomplete report prose about eigenfunction limits and a conditional scarring result, not a resolved standalone problem. The assessment is retained as dated provenance; it is not a new comprehensive literature or open-status certification.
\subsection{Short English statement}
Determine whether some high-energy eigenfunctions of an irrational point-scatterer torus concentrate in momentum directions.
\subsection{Sources}
\begin{itemize}
\item[\hypertarget{source:statement:1}{S1}] ULAM.AI and the UnsolvedMath Contributors, UnsolvedMath, catalogue entry. Catalogue curation under CC BY 4.0; definitions and English exposition added locally.
\url{https://huggingface.co/datasets/ulamai/UnsolvedMath}.
\item[\hypertarget{source:statement:2}{S2}] The Analytic Theory of Automorphic Forms (2011), Oberwolfach Reports, DOI 10.4171/owr/2011/43, pp. 2397--2458; official report source and mathematical context.
\url{https://doi.org/10.4171/owr/2011/43}.
\item[\hypertarget{source:statement:3}{S3}] Official EMS Press PDF of the same Oberwolfach report; PDF pages 41--42, printed pages 2437--2438.
\url{https://ems.press/content/serial-article-files/46359}.
\end{itemize}
\label{end:statement}
\end{document}
